% ===== 第 1 章 c01b：Chan 2014，原文 pp.6--9 =====
% [原文 p.6]
遗憾的是，在写作本文时，五次三维流形上是否存在特殊拉格朗日环面纤维化仍属未知。

非紧的特殊拉格朗日纤维化的例子则容易得多。Harvey 与 Lawson 关于校准几何的著名论文 [86] 给出了其中最简单的例子：由
\[
  f : \C^3 \to \R^3,\qquad
  (z_1,z_2,z_3) \mapsto \bigl(\operatorname{Im}(z_1 z_2 z_3),\, |z_1|^2-|z_2|^2,\, |z_1|^2-|z_3|^2\bigr),
\]
定义的映射是一个特殊拉格朗日纤维化，其纤维在 $\C^3$ 的对角 $T^2$-作用下不变。这个例子后来被 Goldstein [62] 与 Gross [69] 的独立工作大幅推广。他们通过保持自然全纯体积形式的相容 $T^{n-1}$-作用，在任一环面 Calabi--Yau $n$-重（必然非紧）上构造了显式的特殊拉格朗日环面纤维化。这些例子的判别轨迹具有实余维二，并且可以显式描述。

另一组对 SYZ 镜像对称与特殊拉格朗日几何的发展具有历史影响的非紧例子，由 Joyce [95] 发现。人们曾经相信特殊拉格朗日纤维化总是光滑的，从而判别轨迹具有余维二；但 Joyce 的例子表明这不太可能是对的。他构造了显式的 $S^1$-不变特殊拉格朗日纤维化，它们只是\emph{分段光滑}的，且判别轨迹具有\emph{实余维一}。这类纤维化的奇点集是一个 Riemann 曲面，其 amoeba 形状的像给出余维一的判别轨迹。Joyce 还论证说，他的例子展示了特殊拉格朗日纤维化判别轨迹的\emph{一般}行为。

Joyce 这一开创性工作大大加深了我们对特殊拉格朗日纤维化奇点的理解，同时也迫使我们重新思考 SYZ 猜想的表述。最初，我们期望一对镜像 Calabi--Yau 流形 $X$ 与 $\check X$ 应当有到同一底空间 $B$ 的特殊拉格朗日环面纤维化，使得它们的判别轨迹重合。但 Joyce 的例子表明：一方的判别轨迹可以是余维一的，而另一方的判别轨迹可以是余维二的。人们所能期望的最好情形是：$X$ 与 $\check X$ 都容许到同一底空间 $B$ 的特殊拉格朗日环面纤维化，并且当双方都趋于大复结构极限时，这些纤维化的判别轨迹收敛到同一个余维二的极限。

更准确地说，设 $\mathcal{X} \to D$ 与 $\check{\mathcal{X}} \to D$ 是互为镜像的 Calabi--Yau 流形的极大幂零退化，其中 $D$ 是单位圆盘，$0 \in D$ 对应于镜像对的大复结构极限。我们取一列 $\{t_i\} \subset D$ 收敛到 $0$，并设 $g_i$ 与 $\check g_i$ 分别是 $\mathcal{X}_{t_i}$ 与 $\check{\mathcal{X}}_{t_i}$ 上的 Ricci 平坦度量，它们被归一化到具有固定直径 $C$。那么我们期望：
\begin{enumerate}
  \item $(\mathcal{X}_{t_i},g_i)$ 与 $(\check X_{t_i},\check g_i)$ 存在收敛子列，分别（在 Gromov--Hausdorff 意义下）收敛到度量空间 $(B_\infty,d_\infty)$ 与 $(\check B_\infty,\check d_\infty)$；
  \item 空间 $B_\infty$ 与 $\check B_\infty$ 是带奇点的仿射流形，且都同胚于 $S^n$；
% [原文 p.7]
  \item 在某个实余维二的轨迹 $\Gamma \subset B_\infty$（相应地，$\check\Gamma \subset \check B_\infty$）之外，$d_\infty$（相应地，$\check d_\infty$）由某个 Monge--Amp\`ere 度量诱导；并且
  \item Monge--Amp\`ere 流形 $B_\infty \setminus \Gamma$ 与 $\check B_\infty \setminus \check\Gamma$ 互为 Legendre 对偶。
\end{enumerate}

SYZ 猜想的这个\emph{极限版本}由 Gross--Wilson [84] 与 Kontsevich--Soibelman [106] 独立提出。事实上，理解 Ricci 平坦度量极限行为的一般问题，是 Yau 在其著名的公开问题清单 [171, 172] 中提出的。受镜像对称的 SYZ 图景启发，过去 15 年间这一问题得到了广泛研究，Gross--Wilson [84]、Tosatti [158, 159]、Ruan--Zhang [145]、Zhang [173]、Rong--Zhang [141, 140] 以及更近期的 Gross--Tosatti--Zhang [83, 82] 都取得了实质性进展。

度量空间 $B_\infty$ 与 $\check B_\infty$ 应当被看作 Calabi--Yau 流形族上 SYZ 纤维化底空间的极限。应用上述图景，人们可以如下尝试构造 Calabi--Yau 流形极大幂零退化 $\mathcal{X} \to D$ 的镜像：首先确定 Gromov--Hausdorff 极限 $B_\infty$；然后取 $B_\infty \setminus \Gamma$ 的 Legendre 对偶 $\check B_0$，并尝试紧化商 $\check X_0 := T\check B_0/\Lambda$ 以得到正确的镜像。不幸的是，这一朴素的途径行不通：由于仿射结构在 $\check B_\infty$ 的奇点周围具有非平凡单值群，$\check X_0$ 上的自然复结构并非整体定义的。为得到修正后的镜像，需要考虑全纯圆盘的贡献，从而形变 $\check X_0$ 上的复结构。

\section{紧 Calabi--Yau 流形的 SYZ}
人们期望：利用半平坦 SYZ 镜像对称得到的辛结构可以自然地紧化，从而在镜像上给出整体的辛结构。确实，Casta\~no-Bernard 与 Matessi 的工作 [15] 表明，Gross 在 [70] 中构造的拓扑 Calabi--Yau 紧化可以做成辛紧化，由此得到一对紧辛六维流形；它们同胚于已知的 Calabi--Yau 三维流形的镜像对（例如五次三维流形及其镜像），并且带有拉格朗日环面纤维化，其底空间是互为 Legendre 对偶的带奇点整仿射流形。

另一方面，正如我们之前提到的，原始的 SYZ 方案 [154] 已经预期：镜像上的 Ricci 平坦度量应当以来自全纯圆盘的瞬子修正而不同于半平坦 Calabi--Yau 度量 [65]；这些全纯圆盘的边界绕 SYZ 纤维化纤维中的非平凡 1-循环。由于镜像上的度量由其辛结构与复结构唯一确定，自然地期望所有这些瞬子修正都贡献于镜像上复结构的扰动。这确实正是 SYZ 猜想背后的关键思想。由于全纯圆盘可以粘接成全纯曲线，这就解释了为什么镜像对称能够给出计数预言。

现在给定一个带奇点的仿射流形 $B$。设 $\Gamma \subset B$ 是奇异轨迹，记 $B_0 = B \setminus \Gamma$ 为光滑部分。那么人们想构造一个复流形，它是一个小形变 $X_0 := TB_0/\Lambda$ 的紧化。这称为\emph{重构问题}，它正是 Gross 与 Siebert 的代数几何 SYZ 纲领 [76, 77, 78, 79] 的核心所在。

% [原文 p.8]
重构问题最早由 Fukaya [50] 着手处理。他试图通过直接求解控制 $X_0$ 上复结构形变的 Maurer--Cartan 方程来找到合适的扰动。在维数为二的情形，他的启发式论证表明：扰动应当来自 $B$ 中的梯度流树，这些树的端点从奇异集 $\Gamma$ 发出；后者应当是当人们趋近大复结构极限时，界定拉格朗日纤维与 SYZ 纤维化奇点的全纯圆盘的极限。Fukaya 提出了一系列漂亮的猜想，解释量子修正如何贡献于镜像上的复结构，并给出了直观清晰的图景；但遗憾的是，使他的论证严格化所需的分析似乎遥不可及。

Kontsevich 与 Soibelman [107] 通过使用刚性解析空间绕开了 Fukaya 论证中的分析困难。他们从 $S^2$ 上一个带有 24 个奇点、且每个奇点周围仿射结构的单值群都是最简形式
\[
  \begin{pmatrix} 1 & 1 \\ 0 & 1 \end{pmatrix}
\]
的整仿射结构出发，成功构造了一个非阿基米德解析 K3 曲面。基本想法是：对 Fukaya 构造中的每条梯度流线附加一个自同构，并用这些自同构修改镜像中复坐标卡之间的粘接，从而消解因判别轨迹周围仿射结构的非平凡单值群而产生的坐标卡不相容性。他们论证中的一个关键步骤是一个关键引理：当两条梯度流线相交时，总可以加入新的流线并附加新的自同构，使得绕每个交点的复合是恒等映射。这称为\emph{散射现象}，它保证沿一条路径所穿过的各条流线所附自同构的复合与路径选择无关，从而保证修改后的粘接是相容的。

大约在同一时期，Gross 与 Siebert 启动了他们引人注目的纲领 [76, 77, 78, 79]，旨在为 SYZ 猜想建立代数几何的途径。受上一节讨论的 SYZ 猜想极限版本以及 Kontsevich 的一个观察（Gromov--Hausdorff 极限大致就是退化的对偶相交复形）的启发，他们表述了一个代数几何的 SYZ 程序来构造镜像。更详细地说，从 Calabi--Yau 流形的环面退化出发，第一步是构造对偶相交复形；然后取（离散）Legendre 变换，并尝试从 Legendre 对偶重构镜像的环面退化。这样一来，他们可以完全忘掉 SYZ 纤维化。其主张是：所有信息都编码在对偶相交复形的热带几何中；后者是一个带奇点的整仿射流形，扮演 SYZ 纤维化底空间的角色。

利用 Kontsevich 与 Soibelman 的上述关键引理，并辅以采用对数结构与热带几何技巧等许多新想法，Gross 与 Siebert 最终成功地给出了\emph{任意}维数下重构问题的一个解 [79]。更准确地说，给定任意满足某些假设、并带有若干附加结构（如多面体分解）的带奇点整仿射流形，他们构造出 Calabi--Yau 流形的环面退化，并且可以用 $B$ 中的热带树显式且典范地描述它。此外，他们构造的 Calabi--Yau 流形定义在 $\C$ 上，而不是刚性解析空间。

% [原文 p.9]
现在的目标是通过穿越热带世界来获得对镜像对称的概念性理解。在 B-侧，Gross 与 Siebert 猜想：他们构造中的形变参数是一个典范坐标，并且 Calabi--Yau 流形族的周期积分可以用 $B$ 中的热带圆盘来表达。他们已在局部情形（如 [79, 注 5.1] 中的局部 $\PP^2$ 例子）给出了一些证据，并正在处理一般情形。

在 A-侧，人们希望完全在环面退化的中心纤维上工作，来理解光滑纤维的 Gromov--Witten 理论；而该中心纤维的对偶相交复形正是 Gross 与 Siebert 出发时所用的仿射流形。Abramovich--Chen [31, 4] 与 Gross--Siebert [81] 最近的独立工作发展了 对数 Gromov--Witten 不变量理论，推广了 A.-M. Li 与 Ruan [114]、Ionel 与 Parker [93, 94] 以及 Jun Li [115] 关于相对 Gromov--Witten 理论的工作，朝着这一目标迈出了重要一步。若能进一步证明热带曲线/圆盘与全纯曲线/圆盘之间的一般对应定理（沿 Mikhalkin [129, 130]、Nishinou--Siebert [138] 与 Nishinou [137, 136] 工作的路线），我们就能把 Calabi--Yau 流形的 A-侧（即 Gromov--Witten 理论）与热带世界联系起来。

尽管还有许多工作要做，这已勾勒出一幅令人满意的图景，用热带几何来解释镜像对称的几何。关于 Gross--Siebert 纲领的概览，我们请读者参阅其创立者们的出色综述 [80, 66]。

\section{非紧 Calabi--Yau 流形的 SYZ}
特殊拉格朗日环面纤维化例子的缺乏，是对紧 Calabi--Yau 流形实施原始 SYZ 方案的一个主要障碍（也许这也是 Gross 与 Siebert 想要发展代数几何版本的原因之一）。但在非紧情形有大量例子，可以显式地找到特殊拉格朗日环面纤维化，例如 Goldstein [62] 与 Gross [69] 在环面 Calabi--Yau 流形情形构造的那些。此外，按 Fukaya、Oh、Ohta 与 Ono 的工作 [55, 56, 53]，在环面情形下，计数从开 Riemann 曲面到该流形的映射的开 Gromov--Witten 不变量是良定义的（甚至 $S^1$-等变情形也可定义，如 Liu 的学位论文 [121] 所做）。因此，对环面 Calabi--Yau 流形直接实施 SYZ 方案而不退回到热带世界，是完全合理的。

这就把我们带到局部镜像对称的领域。局部镜像对称最初是镜像对称技巧对紧 Calabi--Yau 三维流形中 Fano 曲面的应用，可以通过在 K\"ahler 模空间与复模空间中取某些极限，从紧 Calabi--Yau 流形的镜像对称用物理论证导出 [99]。由于这种镜像对称提供了大量有趣的例子与预言，它一直吸引着物理学家与数学家的广泛关注 [112, 35, 89, 69, 70, 155, 101, 63, 91, 92, 47, 48, 102, 146]。

设 $X$ 是一个 $n$ 维环面 Calabi--Yau 流形，它必然非紧。为实施 SYZ 构造，我们考虑 Goldstein 与 Gross 构造的特殊拉格朗日环面纤维化 $\mu : X \to B$；这种纤维化是非环面的，意即它不是 $X$ 上 Hamiltonian $T^n$-作用所对应的通常矩映射。该纤维化的判别轨迹已由 Gross 详细分析，并可显式描述。
